% Saran penelitian selanjutnya
\section{Rekomendasi untuk Penelitian Selanjutnya}

Penelitian mendatang perlu mengupayakan permohonan data agregat tanpa identitas peserta (seperti matriks korelasi antarsubtes atau parameter skor komposit nasional) kepada Panitia SNPMB melalui mekanisme keterbukaan informasi publik (UU No. 14 Tahun 2008) atau skema kolaborasi riset institusional. Data mikro resmi memungkinkan analisis tanpa data sekunder, sedangkan parameter resmi tingkat populasi mempercepat penelitian, dan salah satu maupun keduanya bermanfaat. Standardisasi komputasi terbuka melengkapi langkah itu. Adopsi arsitektur \textit{pipeline} komputasi terstandardisasi (mengikuti pedoman Project TIER berbasis skrip R terbuka) memfasilitasi replikasi independen dan audit ilmiah terbuka pada studi psikometri pendidikan di Indonesia. Paket replikasi R terbuka pada paper ini mempermudah pembaruan kalkulasi apabila data mikro resmi, parameter resmi tingkat populasi, atau keduanya dirilis secara terbuka.

Kolaborasi dengan lembaga uji coba (\textit{try-out}) atau platform bimbingan belajar berskala nasional berpotensi menghimpun sampel berstrata yang lebih mencerminkan spektrum kemampuan populasi. Stratifikasi juga dapat memakai proksi pilihan program studi untuk membandingkan profil skor peserta bidang Saintek dan Soshum.

Replikasi pada UTBK-SNBT tahun lain dan pada ujian seleksi serupa di Indonesia diperlukan untuk menilai apakah estimasi korelasi dan struktur faktor yang diperoleh bersifat khas 2026, termasuk akibat format dua skor LBI.

Riset berbasis \textit{crowdsourcing} berikutnya disarankan mengintegrasikan algoritma deteksi citra berbasis kecerdasan buatan untuk mengidentifikasi anomali dan pemalsuan sertifikat secara objektif dan sistematis.

Langkah pertama adalah pengujian diferensiasi kemampuan (SLODR), yaitu apakah korelasi antarsubtes berubah menurut tingkat kemampuan. Pengujian itu menuntut data mikro populasi atau sampel yang mencakup rentang skor penuh. Data semacam itu dapat dianalisis melalui model faktor multigrup atau analisis faktor yang membiarkan muatan berubah menurut tingkat kemampuan (\textit{moderated factor analysis}) dan dibandingkan dengan koreksi Case III yang mengasumsikan korelasi tunggal. Simulasi skenario pada paper ini dapat menjadi acuan awal, sedangkan kemiringan SLODR yang sebenarnya hanya dapat ditaksir dari data empiris itu. Langkah kedua adalah mengembangkan varian koreksi \textit{restriction of range} yang mengakomodasi struktur korelasi berjenjang (mengikuti struktur dua faktor Literasi dan Numerasi hasil CFA) untuk menggantikan asumsi korelasi seragam tunggal antarpasangan subtes. Langkah ini relevan karena rasio $SD$ tersirat bervariasi antarsubtes pada paper ini (Tabel~\ref{tab:diagnostik}).

Pemanfaatan data mikro populasi di masa depan memungkinkan pengujian langsung terhadap bentuk distribusi skor (kemiringan, kurtosis, dan kecocokan transformasi normalisasi seperti Box-Cox atau Johnson) serta estimasi persentil secara empiris tanpa asumsi normal.